
%(BEGIN_QUESTION)
% Copyright 2015, Tony R. Kuphaldt, released under the Creative Commons Attribution License (v 1.0)
% This means you may do almost anything with this work of mine, so long as you give me proper credit

Les og lag disposisjon av introduksjonen og underkapitlet ``Potential Transformers'' i delen ``Electrical Sensors'' i kapitlet ``Electric Power Measurement og Control'' i læreboken {\it Lessons In Industrial Instrumentation}. Noter sidetall der viktige illustrasjoner, fotografier, ligninger, tabeller og andre relevante detaljer finnes. Forbered deg på en faglig diskusjon med lærer og medstudenter om begrepene og eksemplene i lesingen.

\underbar{file i01243}
%(END_QUESTION)




%(BEGIN_ANSWER)


%(END_ANSWER)





%(BEGIN_NOTES)

En potensialtransformator (spenningstransformator) transformerer ned og isolerer høy linjespenning til et nivå som er trygt for panelinstrumenter som voltmetre. Fast forhold i en PT (eller VT) gir en proporsjonal representasjon av enten fase- eller linjespenning (avhengig av hvordan transformatoren er koblet). 120 VAC er industristandard for 100\% normal linjespenning, og dette muliggjør samvirke mellom ulike fabrikat og modeller av PT-er og instrumenter.

\vskip 10pt

En CCVT (capacitively coupled voltage transformer) brukes i svært høyspente anlegg for å redusere linjespenningen videre ned til nominelt 120 VAC-signalforsyning. Den bruker seriekondensatorer som en spenningsdeler før nedtransformeringstrinnet.







\vskip 20pt \vbox{\hrule \hbox{\strut \vrule{} {\bf Forslag til sokratisk diskusjon} \vrule} \hrule}

\begin{itemize}
\item{} Hva er hensikten med å bruke PT (eller VT) i elektrisk effektinstrumentering?
\item{} Hva er standard utgangssignalområde for en PT?
\item{} Forklar hvorfor PT-er har standardiserte utgangssignalområder.
\item{} I hvilke typer effektsystemer forventer du å se CCVT brukt i stedet for en rent induktiv VT?
\end{itemize}


%INDEX% Leseoppgave: Lessons In Industrial Instrumentation, potensialtransformatorer

%(END_NOTES)
